\section{Metric topology with explicit computational data}
\label{sec:metric}

\subsection{What the metric replacements assert}

The four entries treated here are Tychonoff's theorem,\footnote{\url{https://mathworld.wolfram.com/TychonoffTheorem.html}}
Urysohn's lemma,\footnote{\url{https://mathworld.wolfram.com/UrysohnsLemma.html}}
Urysohn's metrization theorem,\footnote{\url{https://mathworld.wolfram.com/UrysohnsMetrizationTheorem.html}}
and the compact-open topology.\footnote{\url{https://mathworld.wolfram.com/Compact-OpenTopology.html}}  They concern
infinite objects in their usual formulations.  The replacements below are
families of arithmetic assertions about names, moduli, finite nets, and
finite computations.  The displayed real-valued notation is an explanatory
interpretation of those assertions; it does not add a sort of arbitrary
real numbers or arbitrary subsets to \ipra.

A named real is given by a fixed primitive-recursive rational approximation
procedure with its error estimate.  We write $q(k)\simeq_k r$ for
$|q(k)-r|\leq 2^{-k}$.  A named point has enough data to compute the
rational approximations of the distances used below.  All operations are
understood in a fixed, explicitly given primitive-recursive presentation.
When a family of spaces is indexed by $i$, the procedures involving $i$
must be uniformly primitive recursive.  The phrase ``given a procedure''
means a particular function symbol with its defining equations and the
stated hypotheses, or a proof schema after substituting that procedure.
It never means that arbitrary primitive-recursive program codes can be
uniformly evaluated by another primitive-recursive program.

For a metric space $X$, the following data have different roles.
\begin{itemize}
\item A modulus for a map $f:X\to Y$ is a function $\mu$ satisfying
\begin{equation}\label{eq:metric-modulus}
d_X(x,y)\leq 2^{-\mu(k)}\quad\Longrightarrow\quad
 d_Y(f(x),f(y))\leq 2^{-k}.
\end{equation}
Replacing $\mu(k)$ by $\max_{j\leq k}\mu(j)$ makes it monotone.
\item A net is a finite list $P_k=(p_{k,j})_{j<M(k)}$ together with a
selector $\sigma_k(x)<M(k)$ and a proof of
$d_X(x,p_{k,\sigma_k(x)})\leq 2^{-k}$.  The selector is useful data;
it is not obtained by deciding inequalities between arbitrary reals.
\item A Cauchy modulus for a sequence $(x_n)$ is a function $m$ such that
$n,l\geq m(k)$ implies $d_X(x_n,x_l)\leq 2^{-k}$.
\item A completion operation, when assumed, takes the particular supplied
Cauchy data to a point name and a convergence modulus.  The existence of
an arbitrary limit is not being used as a choice principle.
\end{itemize}
Distances can be truncated at $1$ without changing the metric topology:
$\bar d=\min(1,d)$ is a metric because
$\min(1,a+b)\leq\min(1,a)+\min(1,b)$.  The triangle inequality for
$\bar d$ follows immediately.  The balls of radius less than $1$ coincide
with the corresponding $d$-balls.  Finite minima and maxima of named reals
are computed by $\max(a,b)=(a+b+|a-b|)/2$, so their use does not invoke
real-order decidability.

\begin{remark}[Finite-demand meaning]\label{rem:metric-finite-demand}
Every precision-$k$ conclusion below depends on explicitly bounded lists,
finite sums, and rational inequalities.  Rational inequalities are tested
by cross multiplication with positive denominators.  A proof can replace
all finite bounded quantifiers by their primitive-recursive Boolean folds.
The correctness invariant for the relevant loop is then quantifier free
with the input codes as parameters.  Analytic notation abbreviates these
uniform estimates, not an unproved global soundness theorem for all such
notations.  The Lean file supplied with this chapter checks some abstract
metric-interface lemmas and states the remaining implementation targets;
it is not a derivation of every displayed theorem inside an encoded
\ipra\ proof calculus.
\end{remark}

\subsection{Products: a constructive Tychonoff replacement}

\begin{theorem}[Finite products of explicit nets]\label{thm:finite-metric-product}
Let $X_0,\ldots,X_{r-1}$ have nets at every precision, and equip their
product with $d_{\max}(x,y)=\max_{i<r}d_i(x_i,y_i)$.  The Cartesian
product of the precision-$k$ nets is a precision-$k$ net, of size
$\prod_{i<r}M_i(k)$.  Its selector and list are primitive recursive.
For $r=0$ the product consists of one empty tuple.
\end{theorem}
\begin{proof}
For each coordinate choose the supplied index $\sigma_{i,k}(x_i)$.
Every coordinate distance to the resulting tuple is at most $2^{-k}$,
so their maximum is at most $2^{-k}$.  Mixed-radix enumeration of the
finite tuple of indices gives both the list and the selected index.
Only bounded products and bounded loops are used.  If a factor is empty,
there is no input product point to select; its empty net is appropriate.
\end{proof}

For a countable family assume bounded metrics $0\leq d_i\leq 1$ and
specified base points $p_i$.  A product point is a uniformly presented
family $(x_i)$, not an arbitrary externally chosen sequence.  Define
\begin{equation}\label{eq:product-metric}
 D(x,y)=\sum_{i=0}^{\infty}2^{-i-1}d_i(x_i,y_i),\qquad
 D_N(x,y)=\sum_{i<N}2^{-i-1}d_i(x_i,y_i).
\end{equation}
The sum is shorthand for a name with the certified bound
\begin{equation}\label{eq:product-tail}
0\leq D(x,y)-D_N(x,y)\leq 2^{-N}.
\end{equation}

\begin{theorem}[Weighted product metric and its neighborhoods]
\label{thm:weighted-product}
The procedure in \eqref{eq:product-metric} defines a named metric.
It satisfies
\begin{equation}\label{eq:coordinate-bound}
 2^{-i-1}d_i(x_i,y_i)\leq D(x,y).
\end{equation}
Its neighborhoods agree with the finite-coordinate product neighborhoods,
with explicit transformations in both directions.
\end{theorem}
\begin{proof}
The finite geometric sum gives
$\sum_{i=N}^{L-1}2^{-i-1}\leq 2^{-N}$ for every $L\geq N$;
this proves the Cauchy estimate defining the sum.  For a precision-$k$
name, sum the first $k+2$ terms using distance approximations of error
$2^{-k-2}$.  Their weighted error is at most $2^{-k-2}$, and the omitted
tail has the same bound.  Thus the total error is at most $2^{-k-1}$.
Nonnegativity, symmetry, and the triangle inequality hold first for finite
sums and then at every demanded precision.  If $D(x,y)=0$, the nonnegative
summand bound gives $d_i(x_i,y_i)=0$ for every fixed $i$, hence coordinate
identity.  No decision between equality and inequality is involved.

For a finite cylinder requiring $d_i(x_i,y_i)<r_i$ for $i<N$, the positive
rational radius
$\rho=\min_{i<N}2^{-i-1}r_i$ satisfies
$D(x,y)<\rho\Rightarrow d_i(x_i,y_i)<r_i$ for all its coordinates.
The empty cylinder is the whole space.  Conversely, to obtain
$D(x,y)<2^{-k}$, put $N=k+2$ and require
$d_i(x_i,y_i)<2^{-k-2}$ for $i<N$.  The partial sum is less than
$2^{-k-2}$ and the tail is at most $2^{-k-2}$; their sum is less than
$2^{-k-1}<2^{-k}$.  General positive rational radii are reduced to dyadic
ones by a bounded calculation on their numerator and denominator.
\end{proof}

\begin{theorem}[Countable products of explicit nets]
\label{thm:countable-net-product}
Suppose the base points, nets, and selectors are supplied uniformly in
$i$.  The product has a primitive-recursive net of size
\begin{equation}\label{eq:product-net-size}
 M(k)=\prod_{i<k+2}M_i(k+2).
\end{equation}
A net point uses a net coordinate in positions $i<k+2$ and the fixed base
point $p_i$ in every remaining position.  Its distance from the selected
input point is at most $2^{-k-1}$.
\end{theorem}
\begin{proof}
For the first $k+2$ coordinates the error is at most $2^{-k-2}$.
Their weighted contribution is at most that quantity.  The unused tail
is at most $2^{-k-2}$.  Sum these bounds.  Listing and selecting the finite
coordinate tuple uses the algorithm of
Theorem~\ref{thm:finite-metric-product}; evaluating its other coordinates
uses the given base-point procedure.  In particular, no countable choice
of points in unspecified nonempty factors occurs.
\end{proof}

\begin{theorem}[Transport of supplied completion procedures]
\label{thm:product-completion}
Suppose coordinate completion is available uniformly: from the specific
coordinate Cauchy data it produces a name $z_i$ and a convergence modulus
$c_i$.  Then completion of the corresponding product sequence is
primitive recursive.  A valid product convergence modulus is
\begin{equation}\label{eq:product-completion-rate}
 c(k)=\max_{i<k+2}c_i(k+2).
\end{equation}
Product Cauchy data with modulus $m$ furnish coordinate-$i$ Cauchy data
with modulus $k\mapsto m(k+i+1)$.
\end{theorem}
\begin{proof}
The second assertion follows from
$d_i\leq 2^{i+1}D$.  Apply the given coordinate completion procedure to
these data, uniformly in $i$.  If $n\geq c(k)$, each of the first $k+2$
coordinate distances from $z_i$ is at most $2^{-k-2}$.  The same
partial-sum and tail estimate used for the net gives
$D(x_n,z)\leq 2^{-k-1}$.  Each finite request to the product name makes
only finitely many calls to the supplied uniform procedure.
\end{proof}

Theorems~\ref{thm:countable-net-product} and
\ref{thm:product-completion} are a useful version of countable product
compactness when compactness is represented by completeness data and
finite nets.  They do not assert the open-cover compactness of an
arbitrary product or extract a modulus from an arbitrary cover.
The distinction is essential: the MathWorld Tychonoff entry discusses an
arbitrary family, and explicitly records its relation to choice.
There is also a classical edge case in the entry's opening ``if and only
if'': an empty product space may be compact even if one of its other
factors is noncompact.  The converse from compact product to compact
factors requires a nonempty product, for example specified base points.

\begin{proposition}[Uniform images preserve nets]\label{prop:image-nets}
If $f$ has modulus $\mu$ and $P_k$ is a net for $X$, then
$(f(p))_{p\in P_{\mu(k)}}$ covers the represented image of $f$ at
precision $k$, with a selector supplied by the preimage presentation.
\end{proposition}
\begin{proof}
For an image point presented as $f(x)$, use
$f(p_{\mu(k),\sigma_{\mu(k)}(x)})$ and apply
\eqref{eq:metric-modulus}.  This does not choose a preimage from a bare
extensional membership assertion.  Completeness of an image is not
claimed merely from this net calculation.
\end{proof}

\subsection{Urysohn functions from located distance data}

A distance function to a general closed set need not be computable from
a negative description of that set.  We therefore include the functions
needed for the construction in the input.  Let $a,b:X\to[0,\infty)$ be
supplied $1$-Lipschitz named functions, with $a=0$ on $A$ and $b=0$ on
$B$.  They may be the located distances to $A$ and $B$.

\begin{theorem}[Urysohn separation with a positive gap]
\label{thm:metric-urysohn}
If $a(x)+b(x)\geq\delta>0$ for all $x$, where $\delta$ is a positive
rational input, then
\begin{equation}\label{eq:urysohn-function}
 u(x)=\frac{a(x)}{a(x)+b(x)}
\end{equation}
is a named function with values in $[0,1]$, is $1/\delta$-Lipschitz,
and equals $0$ on $A$ and $1$ on $B$.  Its approximation procedure and
modulus are primitive recursive in the supplied data.
\end{theorem}
\begin{proof}
Write $a=a(x)$, $b=b(x)$, $a'=a(y)$, $b'=b(y)$, and $s=a+b$,
$s'=a'+b'$.  The denominator is positive and
\[
 \left|\frac{a}{s}-\frac{a'}{s'}\right|
 =\frac{|ab'-a'b|}{ss'}
 \leq\frac{a|b'-b|+b|a-a'|}{ss'}
 \leq\frac{d(x,y)}{s'}\leq\frac{d(x,y)}{\delta}.
\]
This proof gives the sharper constant $1/\delta$, with no factor of $2$.
The interval bound follows from $0\leq a\leq s$.  On $A$ the numerator
vanishes, while on $B$ one has $s=a\geq\delta$.

For computation approximate $a,b$ with error $\eta$ and clamp the
approximations to be nonnegative.  Clamping cannot increase the error.
If $2\eta\leq\delta/2$, their approximate sum $\widehat s$ is at least
$\delta/2$.  The same numerator calculation gives
\[
 \left|\frac{a}{s}-\frac{\widehat a}{\widehat s}\right|
 \leq\frac{\eta}{\widehat s}\leq\frac{2\eta}{\delta}.
\]
Choose a natural number $h$ with
$2^{-h}\leq\min(1,\delta)/8$, and evaluate at precision $k+h$.
The error is at most $2^{-k-2}$.  Such an $h$ can be obtained by a
bounded search using the denominator of the positive rational $\delta$.
A uniform-continuity modulus is obtained by choosing $\mu(k)$ so that
$2^{-\mu(k)}\leq\delta 2^{-k}$.  Every division is by an explicitly
positive rational approximation.
\end{proof}

\begin{proposition}[Local-gap Urysohn separation]
\label{prop:local-urysohn}
A global lower bound is unnecessary for pointwise continuity.  It is
enough to supply, for every presented $x$, a positive rational
$\delta_x\leq a(x)+b(x)$, by a primitive-recursive procedure.  Formula
\eqref{eq:urysohn-function} then has a computable local modulus: if
$d(x,y)\leq\delta_x/4$, then
\[
 |u(x)-u(y)|\leq 2d(x,y)/\delta_x.
\]
\end{proposition}
\begin{proof}
Since $a+b$ is $2$-Lipschitz,
$a(y)+b(y)\geq\delta_x-2d(x,y)\geq\delta_x/2$.
Use the numerator estimate from the preceding proof with this denominator.
Evaluation at $x$ uses $\delta_x$ exactly as in that proof.  Thus a
local modulus may demand simultaneously
$d(x,y)\leq\delta_x/4$ and
$d(x,y)\leq\delta_x 2^{-k-1}$.
Disjointness alone does not supply the distance functions or this local
lower-bound procedure.
\end{proof}

\begin{corollary}[Finite rational metric sets]
\label{cor:finite-metric-urysohn}
Let $A,B$ be nonempty disjoint finite subsets of a rational metric space,
and suppose their displayed points are pairwise distinguishable by the
metric.  Put
\[
 a(x)=\min_{p\in A}d(x,p),\quad b(x)=\min_{q\in B}d(x,q),\quad
 \delta=\min_{p\in A,q\in B}d(p,q)>0.
\]
The data satisfy Theorem~\ref{thm:metric-urysohn}.
\end{corollary}
\begin{proof}
Each distance-to-a-point function is $1$-Lipschitz; taking a finite minimum
preserves that bound.  A rational finite minimum attains its value by
bounded comparison.  At minimizing $p,q$, the triangle inequality gives
$\delta\leq d(p,q)\leq d(p,x)+d(x,q)=a(x)+b(x)$.
Finite disjointness and positivity of rational nonzero distances make
the displayed minimum positive.  All lists and comparisons are bounded.
If one of the sets is empty, a constant function supplies the required
nonempty-side boundary value instead.
\end{proof}

\subsection{Metrization by separating coordinates}

The existence of many separating functions does not by itself provide
a computable inverse for the map they define.  The following formulation
records exactly the extra local information that is used.

\begin{theorem}[Coordinate metrization with a neighborhood certificate]
\label{thm:coordinate-metrization}
Let $X$ have a represented neighborhood basis.  Suppose a uniformly
presented family $u_i:X\to[0,1]$ is given with local continuity moduli.
Assume two further pieces of data.
\begin{enumerate}
\item Coordinate equality identifies points: if $u_i(x)=u_i(y)$ for
all $i$, then $x=y$ in the chosen representation.
\item Given a certified basic neighborhood $V$ of $x$, a procedure returns
$N$ and a positive rational $\varepsilon$ such that
\[
 (\forall i<N)\ |u_i(x)-u_i(y)|<\varepsilon\quad\Longrightarrow\quad y\in V.
\]
\end{enumerate}
Then
\[
 d_u(x,y)=\sum_{i\geq0}2^{-i-1}|u_i(x)-u_i(y)|
\]
is a compatible named metric, and $x\mapsto(u_i(x))_i$ is a homeomorphism
onto its represented image in the Hilbert cube.  The inverse
neighborhood transformation can use radius $2^{-N}\varepsilon$.
\end{theorem}
\begin{proof}
The metric axioms and computation follow from
Theorem~\ref{thm:weighted-product}.  Its zero-distance implication is the
first hypothesis.  For continuity at precision $k$, use $N=k+2$ and
intersect the finitely many neighborhoods on which the first $N$
coordinates change by less than $2^{-k-2}$.  The finite sum and tail
give a change less than $2^{-k}$.  The finite intersection and its
refinement are part of the presented basis operations.

Conversely, if $d_u(x,y)<2^{-N}\varepsilon$, then for $i<N$,
\[
 |u_i(x)-u_i(y)|\leq 2^{i+1}d_u(x,y)
 <2^{i+1-N}\varepsilon\leq\varepsilon.
\]
The second hypothesis now gives $y\in V$.  Both neighborhood transformations
are explicit finite calculations.  An inverse is considered only on the
represented image, whose points carry an $X$-name; no algorithm extracting
an arbitrary preimage is implicit.
\end{proof}

This theorem isolates a constructive part of Urysohn's metrization proof.
It does not produce the coordinate functions and their neighborhood
certificates from bare regularity and a bare countable basis.  Those are
additional proof obligations for a particular presentation.

\begin{theorem}[An explicit Hilbert-cube embedding for a separable metric
presentation]\label{thm:dense-hilbert-embedding}
Let $X$ have a specified dense sequence $(p_i)$ and a primitive-recursive
locator $\lambda(x,k)$ satisfying
$d(x,p_{\lambda(x,k)})\leq2^{-k}$.  Define
\[
 u_i(x)=\min(1,d(x,p_i)),\qquad E(x)=(u_i(x))_{i\geq0}.
\]
Then $E$ is an embedding into the Hilbert cube.  Its forward map is
$1$-Lipschitz for the weighted metric, and an explicit inverse local
modulus is
\begin{equation}\label{eq:dense-embedding-inverse}
 D(E(x),E(y))<2^{-(\lambda(x,k+2)+k+2)}
 \quad\Longrightarrow\quad d(x,y)<2^{-k}.
\end{equation}
\end{theorem}
\begin{proof}
Truncation and the reverse triangle inequality give
$|u_i(x)-u_i(y)|\leq\min(1,d(x,y))$; summing the weights gives
$D(E(x),E(y))\leq d(x,y)$.  Put $r=2^{-k}\leq1$ and
$i=\lambda(x,k+2)$.  The hypothesis of
\eqref{eq:dense-embedding-inverse} and the coordinate bound imply
$|u_i(x)-u_i(y)|<r/2$.  Since $u_i(x)\leq r/4$,
$u_i(y)<3r/4<1$.  Thus the truncation is inactive at $y$, and
\[
 d(x,y)\leq d(x,p_i)+d(p_i,y)<r/4+3r/4=r.
\]
This also proves that coordinate equality forces point equality.  Every
index and requested precision in this argument is explicitly computable.
The inverse modulus depends on $x$; no uniform bound on the locator's
index is asserted.
\end{proof}

\subsection{Euclidean examples and explicit homeomorphisms}

The Euclidean-topology and homeomorphism entries supply the classical
terminology for the following concrete constructions.\footnote{\url{https://mathworld.wolfram.com/EuclideanTopology.html};
\url{https://mathworld.wolfram.com/Homeomorphism.html}}

\begin{proposition}[Euclidean norm comparisons]
\label{prop:euclidean-norms}
For $v\in\R^d$, with $d\geq1$,
\[
 \|v\|_\infty\leq\|v\|_2\leq\|v\|_1
 \leq d\|v\|_\infty.
\]
Consequently these metrics induce the same topology by explicit rational
moduli.  Rational grids give nets for bounded rational boxes, and rational
vectors give a dense sequence with a locator for represented points.
\end{proposition}
\begin{proof}
Each $v_i^2$ is at most $\sum_jv_j^2$, so
$|v_i|\leq\|v\|_2$.  Expanding
$(\sum_i|v_i|)^2$ gives $\sum_iv_i^2$ plus nonnegative cross terms,
which proves the middle inequality.  The last inequality sums
$|v_i|\leq\|v\|_\infty$.  Squaring reduces the only square-root
comparisons to rational inequalities with nonnegative sides; the
nonnegative square root has its usual bisection name on a supplied bounded
interval.  For a requested Euclidean error $2^{-k}$, choose a grid of
sup-norm mesh at most $2^{-k}/d$.  Clamping coordinate approximations
to a bounded box preserves the required bounds.  Finite products enumerate
the grid, and standard rational-tuple coding enumerates the dense sequence.
The grid selector uses rational approximations with extra slack, not exact
comparison of a real coordinate with a grid boundary.
\end{proof}

\begin{proposition}[A named homeomorphism between Euclidean space and an
open cube]\label{prop:euclidean-open-cube}
Coordinatewise application of
\[
 h(t)=\frac{t}{1+|t|},\qquad h^{-1}(s)=\frac{s}{1-|s|}
\]
gives a homeomorphism $\R^d\cong(-1,1)^d$.  The forward map is
$1$-Lipschitz in sup norm.  On the subcube $|s_i|\leq1-\delta$,
$\delta>0$, the inverse is $\delta^{-2}$-Lipschitz.
\end{proposition}
\begin{proof}
The two formulas are inverse by elementary substitution, separately on
the nonnegative and nonpositive rational approximations, or directly
using the identities for absolute value.  For two real numbers of the
same sign,
\[
 |h(t)-h(u)|=\frac{|t-u|}{(1+|t|)(1+|u|)}\leq|t-u|.
\]
For opposite signs, $|h(t)-h(u)|=|h(t)|+|h(u)|\leq|t|+|u|=|t-u|$.
The inverse calculation on a same-sign interval has denominator
$(1-|s|)(1-|v|)\geq\delta^2$; across zero use the sum of the two
one-sided estimates.  These inequalities can also be proved by rational
case splits followed by approximation, so they require no trichotomy for
arbitrary named reals.  Evaluation of the inverse at an open-cube point
uses a supplied positive distance from each coordinate to the endpoints.
Taking a finite minimum of those distances supplies $\delta$.
\end{proof}

\subsection{Function spaces and the compact-open topology}

Let $K$ be an inhabited represented metric space equipped with nets.
For maps $f,g:K\to Y$ with supplied moduli, the function
$h(x)=d_Y(f(x),g(x))$ has modulus
\begin{equation}\label{eq:distance-function-modulus}
 \nu(k)=\max\{\mu_f(k+1),\mu_g(k+1)\}.
\end{equation}
Indeed, the reverse triangle inequality gives
$|h(x)-h(y)|\leq d(f(x),f(y))+d(g(x),g(y))$.

\begin{theorem}[Computable uniform metric]
\label{thm:computable-uniform-metric}
With the above data, the quantity
\[
 d_\infty(f,g)=\sup_{x\in K}d_Y(f(x),g(x))
\]
has a primitive-recursive real name.  It is obtained from finite maxima
on nets, and is a metric on maps identified pointwise.  No exact
maximizing point is required.
\end{theorem}
\begin{proof}
At precision $k$ put $\eta=2^{-k-3}$ and take a net at
precision $\nu(k+3)$.  Approximate each of the finitely many values
$h(p_j)$ by $q_j$ with error at most $\eta$, and let $Q_k=\max_jq_j$,
using ordinary rational comparisons.  For every $x$, the selected net
point gives
\[
 h(x)\leq h(p_{\sigma(x)})+\eta\leq Q_k+2\eta.
\]
For a rational maximizing index $j_k$,
$h(p_{j_k})\geq Q_k-\eta$.  Thus $[Q_k-\eta,Q_k+2\eta]$ is a
certified interval for the supremum.  Any two such intervals are
compatible: each upper endpoint bounds every actual value, including the
net value witnessing the other's lower endpoint.  Consequently the
$Q_k$ form a Cauchy name with the displayed error bounds, defining a real
$S$ satisfying $h(x)\leq S$ for every $x$ and
$S\leq Q_k+2\eta$.  If $b$ bounds every $h(x)$ then
$Q_k\leq b+\eta$ for every $k$, hence $S\leq b$ at all precisions.
This proves the claimed supremum property without assuming a general
least-upper-bound axiom.

Symmetry and nonnegativity are inherited pointwise.  For every $x$,
$d(f(x),h(x))\leq d(f(x),g(x))+d(g(x),h(x))$; taking the certified
upper bounds gives the triangle inequality.  If $d_\infty(f,g)=0$,
every pointwise distance is zero, so the maps agree in the chosen
pointwise equivalence.  Every computation made at precision $k$ is a
finite maximum on an explicitly sized net.
\end{proof}

\begin{corollary}[Approximate extrema]
\label{cor:approximate-extrema}
For a real-valued uniformly continuous function $v$ on a represented
inhabited space with nets, a primitive-recursive procedure returns
$x_k\in K$ with
\[
 (\forall x\in K)\quad v(x)\leq v(x_k)+2^{-k}.
\]
An analogous procedure gives an approximate minimum.
\end{corollary}
\begin{proof}
Take a net fine enough that the variation of $v$ between a point and its
selected net point is at most $\eta=2^{-k-2}$.  Approximate the net
values with error at most $\eta$ and choose a largest rational
approximation, at $x_k$.  For every $x$,
\[
 v(x)\leq v(p_{\sigma(x)})+\eta
 \leq q_{\sigma(x)}+2\eta
 \leq q_{x_k}+2\eta
 \leq v(x_k)+3\eta<v(x_k)+2^{-k}.
\]
Apply the same algorithm to $-v$ for minima.  An exact maximizer or a
convergent subsequence of maximizing points is a different assertion and
is not extracted here.
\end{proof}

For the next result take $K=[0,1]^d$ with the sup metric and its usual
dyadic box data.  Write
\[
 [C,U]=\{g:g(C)\subseteq U\}
\]
for a compact-open subbasic condition.  The closed dyadic boxes used
below have explicit product nets and completion procedures.  Thus the
word ``compact'' here comes with the data already constructed, and not
with a universal open-cover elimination principle.

\begin{theorem}[An explicit compact-open neighborhood comparison on a
cube]\label{thm:compact-open-cube}
Let $f:[0,1]^d\to Y$ have a supplied uniform modulus and let
$\varepsilon>0$ be rational.  A primitive-recursive procedure constructs
finitely many closed dyadic boxes $C_j$ covering the cube, points
$p_j\in C_j$, and balls
$U_j=B(f(p_j),\varepsilon/2)$ such that
\begin{equation}\label{eq:compact-open-sandwich}
 B_{\infty}(f,\varepsilon/4)
 \ \subseteq\ \bigcap_j[C_j,U_j]
 \ \subseteq\ B_{\infty}(f,\varepsilon).
\end{equation}
In particular these finitely certified compact-open neighborhoods
and the uniform metric generate the same local neighborhoods for the
represented maps in question.
\end{theorem}
\begin{proof}
Choose a dyadic step $h=2^{-N}$ fine enough that the modulus guarantees
$d_Y(f(x),f(p))\leq\varepsilon/4$ whenever
$\|x-p\|_\infty\leq h$.  Enumerate the grid centers $p_j$ with
coordinates in $\{0,h,2h,\ldots,1\}$ and set
$C_j=\prod_i[\max(0,p_{j,i}-h),\min(1,p_{j,i}+h)]$.
These boxes overlap.  To cover a named point $x$, approximate each
coordinate within $h/4$, clamp the rational approximation to $[0,1]$,
and round it to a nearest grid value.  The resulting center has distance
at most $3h/4$ from $x$, so it selects a containing box with slack.
Enumeration, rational rounding, and the required mesh exponent are
primitive recursive.

If $d_\infty(f,g)<\varepsilon/4$, then for $x\in C_j$,
\[
 d(g(x),f(p_j))
 \leq d(g(x),f(x))+d(f(x),f(p_j))<\varepsilon/2.
\]
This proves the first inclusion.  If $g$ lies in all the displayed
subbasic conditions, select any box containing $x$; then
\[
 d(f(x),g(x))<\varepsilon/4+\varepsilon/2=3\varepsilon/4.
\]
Taking the supremum gives $d_\infty(f,g)\leq3\varepsilon/4<\varepsilon$,
which proves the second inclusion with a uniform positive margin.
The overlap matters constructively: an exact partition into adjacent
closed boxes would ask for a side decision at an arbitrary real boundary.
Our rational-approximation selector makes no such decision.
\end{proof}

For a general compact-open condition $[C,U]$, the analogous finite-demand
statement assumes a supplied margin $\delta>0$ such that
$B(f(x),\delta)\subseteq U$ for all $x\in C$.  Then
$B_\infty(f,\delta)\subseteq[C,U]$ by direct substitution.
No procedure obtaining such a margin from a bare assertion
$f(C)\subseteq U$ has been claimed.  This is the precise limitation of
the comparison with the unrestricted compact-open theorem.

\begin{proposition}[Evaluation with an explicit local modulus]
\label{prop:metric-evaluation}
For modulus-equipped maps on $K$, evaluation $(f,x)\mapsto f(x)$ is
continuous in the uniform metric.  At $(f,x)$, it is enough to require
\[
 d_\infty(f,g)\leq2^{-k-1},\qquad
 d_K(x,y)\leq2^{-\mu_f(k+1)}
\]
to obtain $d_Y(f(x),g(y))\leq2^{-k}$.
\end{proposition}
\begin{proof}
Insert $f(y)$ and use the triangle inequality:
$d(f(x),g(y))\leq d(f(x),f(y))+d(f(y),g(y))
\leq2^{-k-1}+2^{-k-1}$.
The dependence of this local modulus on $f$ is explicit.
\end{proof}

\begin{theorem}[A quantitative exponential-law fragment]
\label{thm:metric-currying}
Suppose $H:K\times Z\to Y$ has a supplied uniform modulus $\mu$ for
the maximum product metric.  Then its curry
$\widehat H(z)(x)=H(x,z)$ has the same modulus as a map to the uniform
metric on $C(K,Y)$.  Conversely, suppose the curry has a modulus
$\alpha$, and all its slices have a common modulus $\beta$ in $x$.
Then $H$ has modulus
\[
 \gamma(k)=\max\{\alpha(k+1),\beta(k+1)\}.
\]
\end{theorem}
\begin{proof}
If $d_Z(z,z')\leq2^{-\mu(k)}$, apply uniform continuity of $H$ to
$(x,z),(x,z')$ for every $x$.  The supremum of their distances is at most
$2^{-k}$.  In the converse direction insert $H(x,z')$:
\[
 d(H(x,z),H(x',z'))
 \leq d(H(x,z),H(x,z'))+d(H(x,z'),H(x',z'))
 \leq2^{-k-1}+2^{-k-1}.
\]
The first term is bounded by the uniform metric between the slices and
the second by their common modulus.  This is a quantitative statement
with supplied uniformity; it is not a derivation of the unrestricted
locally-compact Hausdorff exponential law from bare continuity.
\end{proof}

\subsection{Arithmetic and Lean implementation route}

The finite sum identity behind \eqref{eq:product-tail}, mixed-radix tuple
enumeration, finite maxima, and rational cross-multiplication are the
arithmetic core.  A direct implementation has the following layers.
\begin{enumerate}
\item Encode integers, positive-denominator rationals, lists, tuples, and
finite tables by natural numbers.  Define rational comparison by the
integer cross products, and define every list fold by primitive recursion.
\item Implement each fixed approximation, net, selector, and modulus
procedure as primitive-recursive syntax.  Prove the finite numerical
invariant of each loop by quantifier-free induction.  Rational arithmetic
identities can be proved in the standard definitional extension for
primitive-recursive functions and translated to the selected presentation
of \ipra.
\item Represent a real equation or inequality by its finite precision
tests.  Division carries a positive-denominator certificate; suprema and
infinite sums carry the explicit Cauchy bounds proved above.  The
product construction requests at most $k+2$ coordinates for a
precision-$k$ output.  The embedding requests the one locator index
shown in \eqref{eq:dense-embedding-inverse}.
\item In Lean, first formalize these finite rational statements, then
prove the representation theorems interpreting them as the usual metric
claims.  The supplied \texttt{MetricBlueprint.lean} checks the basic
precision-space product and modulus-composition arguments without
\texttt{sorry}.  Its semantic function fields are interfaces, not claims
that all Lean functions are primitive recursive.
\item For an actual internal endpoint, compile the selected algorithms
and proof invariants into the repository's object-language syntax and
produce its proof predicate.  This last compilation is an explicitly
remaining formalization task; Lean's ordinary kernel checking of an
external semantic theorem does not perform it automatically.
\end{enumerate}
These steps distinguish useful fully proved constructive mathematical
reductions from an assertion that the entire classical topology library
has already been formalized inside \ipra.
